\documentclass[11pt]{article}
\usepackage{mathblatt}
% Präambel des Setzers (werkzeuge/setzer-praeambel.tex, wortgleich); dazu \gteil (Muster T2U)
\makeatletter
\setlength{\mbpfbe}{0mm}% keine Punkte (pruefung.md § 5: nur beim Prüfstein)
\newlength{\gEin}\setlength{\gEin}{0pt}
\let\gSchrift\relax
\renewcommand{\pfkreuzzeile}[1]{\par\noindent\hangindent3.2em\hangafter1\hspace*{1.6em}$\square$\ #1\par}
\newcommand{\gkreuz}[1]{$\square$\ #1\hspace{2em}}
\newcommand{\gkreuzzeile}[2]{\par\noindent\hangindent2.2em\hangafter1\mbox{#1\,$\square$}\ #2\par}% Buchstabe, Option: eine je Zeile (neben Skizze)
% links, was man ansieht, rechts, was man tut; Spaltengrenze überall an derselben Stelle (gemeinsam.md „Satz“)
\newsavebox{\gbox}\newlength{\gLinks}
% v0.10 (Lehrer 13.10.): keine starre Halbierung – die linke Spalte ist so breit wie die Figur (+ 6 mm), die
% Antwort steht direkt daneben; ist die Figur breiter als 0,62 Zeile, steht die Antwort darunter
\newcommand{\gzwei}[2]{\par\smallskip\noindent\sbox{\gbox}{#1}%
  \setlength{\gLinks}{\dimexpr\wd\gbox+6mm\relax}%
  \ifdim\gLinks>0.62\linewidth
    \usebox{\gbox}\par\smallskip\noindent\begin{minipage}{\linewidth}#2\end{minipage}%
  \else
    \begin{minipage}[t]{\gLinks}\vspace{0pt}\usebox{\gbox}\end{minipage}%
    \begin{minipage}[t]{\dimexpr\linewidth-\gLinks\relax}\vspace{0pt}\raggedright #2\end{minipage}%
  \fi\par}
% v0.9 (12.10., pruefung.md § 5 „Marke mit B, Bedingungen rechts“): drittes Argument von pfaufg = Bedingungen
% der Prüfung (v0.10: \pfohnetr = TR durchgestrichen, „FOR“), klein grau im rechten Rand auf der Grundlinie der Nummer, Zeilen mit \\
\newcommand{\gbedrand}[2]{\ifblank{#2}{}{\rlap{\hspace*{\dimexpr#1+2mm\relax}%
  \parbox[t]{15mm}{\raggedright\scriptsize\leavevmode\color{mbgrau}\strut#2}}}}
% Teil eines Bündels (a, b …): Marke links im Rand, Bedingungen rechts im Rand, auf der Zeile des Teils
\newcommand{\gteilrand}[2]{\llap{\mbpfrandmarke{#1}\hspace*{\mbpfnr}}\gbedrand{\linewidth}{#2}}
% Teil eines Bündels: Buchstabe in eigener Spalte, Inhalt eingerückt (\gEin = Einzug, damit die Spaltengrenze von
% \gzwei an derselben Stelle bleibt)
\newenvironment{gbteil}[3]{\par\smallskip\noindent\gteilrand{#2}{#3}\makebox[\mbpfteil][l]{\textbf{#1}}%
  \begin{minipage}[t]{\dimexpr\linewidth-\mbpfteil\relax}\setlength{\gEin}{\mbpfteil}\raggedright\strut\ignorespaces}%
 {\end{minipage}\par}
% Grafik höchstens so breit wie die Spalte (Bündel nebeneinander, Koordinatensystem ohne Karo)
\newsavebox{\gpbox}
\newcommand{\gpassend}[1]{\sbox{\gpbox}{#1}\ifdim\wd\gpbox>\linewidth\typeout{GPASSEND: verkleinert \the\wd\gpbox\space auf \the\linewidth}\resizebox{\linewidth}{!}{\usebox{\gpbox}}\else\usebox{\gpbox}\fi}
% Skizzenfeld (v0.10, Lehrer 13.10.): Karo ab dem linken Rand über die Zeile, „Skizze“ klein grau oben links im Feld
% (wie \pfskizzenfeld der Vorlage); der rechte Rand bleibt den Bedingungen
\newcommand{\gskizze}[1]{\par\smallskip\noindent\begin{tikzpicture}
  \clip (0,0) rectangle ({\linewidth-1mm},{#1*10mm});
  \draw[black!16,line width=0.3pt,step=5mm] (0,0) grid ({\linewidth},{#1*10mm});
  \node[anchor=north west,inner sep=2pt,font=\scriptsize,text=mbgrau,fill=white] at (0,{#1*10mm}) {Skizze};
  \end{tikzpicture}\par}
\RenewDocumentEnvironment{pfaufg}{m m m}{%
  \begin{lrbox}{\mbaufgabebox}\begin{minipage}[t]{\linewidth}%
  \gSchrift\noindent\hspace*{\gEin}\mbpfrandmarke{#1}\gbedrand{\linewidth-\gEin}{#3}%
  \begin{minipage}[t]{\mbpfnr}\strut\textbf{#2}\end{minipage}%
  \begin{minipage}[t]{\dimexpr\linewidth-\gEin-\mbpfnr\relax}\raggedright\strut\ignorespaces}%
 {\end{minipage}%
  \end{minipage}\end{lrbox}%
  \setlength{\mbaufgabehoehe}{\dimexpr\ht\mbaufgabebox+\dp\mbaufgabebox\relax}%
  \ifdefempty{\gkopfnach}{\par\addvspace{5pt}\Needspace*{\mbaufgabehoehe}}{%
    % Stufenkopf, der mit dieser (ersten) Aufgabe zusammenbleibt: Platz für beide verlangen, dann den Kopf setzen
    % (12.10., Quadratische: ein Kopf stand allein unten, die Aufgabe mit Koordinatensystem auf der nächsten Seite)
    \par\Needspace*{\dimexpr\mbaufgabehoehe+40pt\relax}\gkopfnach\gdef\gkopfnach{}\par\addvspace{5pt}}%
  \noindent\usebox{\mbaufgabebox}\par\addvspace{5pt}}
% Stufenkopf vormerken: \gkopfvor{<Kopf>} – die nächste pfaufg setzt ihn, sobald beide zusammen Platz haben
\def\gkopfnach{}
\newcommand{\gkopfvor}[1]{\gdef\gkopfnach{#1}}
% Teilaufgabe mit Skizze links in gzwei: Buchstabe in eigener Spalte, darunter Text und Skizze
\newcommand{\gteil}[2]{\makebox[\mbpfteil][l]{\textbf{#1}}\parbox[t]{\dimexpr0.5\textwidth-\gEin-\mbpfnr-5mm-\mbpfteil\relax}{\raggedright #2}}
% Stufenkopf der Vorlage (mathblatt.sty 2026-10-12: Linie, Abstand, 13 pt, \pfjahre) ohne ihren Vorhalt von
% 14 Zeilen: der Setzer fordert vor dem Kopf 0,16 Seitenhöhe (Kopf + erste Aufgabe), mehr nicht
\setlength{\mbpfstufevorhalt}{0pt}
\makeatother
\newcommand{\karo}[1]{\par\smallskip\noindent\begin{tikzpicture}
  \clip (0,0) rectangle ({\linewidth-1mm},{#1*10mm});
  \draw[black!16,line width=0.3pt,step=5mm] (0,0) grid ({\linewidth},{#1*10mm});
  \end{tikzpicture}\par}
\newcommand{\kk}[1]{$\square$\,#1\hspace{0.9em}}
\newcommand{\linie}[1]{\,{\color{black!45}\rule[-1pt]{#1}{0.4pt}}\,}
\newcommand{\pkt}{\textperiodcentered{}}
% v0.14 (14.10.): Stufenkopf mit fünf Punkten (Option = Jahre für die Bestellkarte) oder „zuletzt 19“
\NewDocumentCommand{\gpunkte}{o m}{#2}
\newcommand{\gzuletzt}[1]{zuletzt #1}
% v0.7 (12.10.): Schülerfrage unter dem Stufenkopf; Ankreuzoption als Bild (Buchstabe und Kästchen darunter)
\newcommand{\gfrage}[1]{\par\nobreak\noindent\hspace*{\mbpfnr}{\small\itshape #1}\par\nobreak\addvspace{2pt}}
\newlength{\gbildbreite}\setlength{\gbildbreite}{0.32\linewidth}% Setzer setzt sie je Reihe (vier Bilder: 0,24)
\newcommand{\gbildoption}[3]{\begin{minipage}[t][#3][b]{\gbildbreite}\centering #2\par\smallskip\mbox{#1\,$\square$}\end{minipage}}
% Lineare (12.10.): Bild mit Buchstabe ohne Kästchen (Diagramm zuordnen); Aussage mit „wahr / falsch“ zum Ankreuzen
\newcommand{\gbildname}[3]{\begin{minipage}[t][#3][b]{\gbildbreite}\centering #2\par\smallskip\mbox{\textbf{#1}}\end{minipage}}
\newcommand{\gwahrfalsch}[2]{\par\noindent\hangindent2.2em\hangafter1\makebox[2.2em][l]{(#1)}#2\hfill\mbox{wahr $\square$\quad falsch $\square$}\par}
\begin{document}

\pfheftstil
\pflaufkopf[Abitur GK \pkt{} Prüfung]{Integral berechnen}{AA140}\pfseitevon
\pfheftkopf{Integral berechnen}{}
\pfstufenzeile{Stammfunktion nachweisen \pkt{} Integral gesucht \pkt{} Grenze oder Parameter gesucht}

% ===================== Stufe: Stammfunktion nachweisen =====================
\gkopfvor{\pfstufekopf{Stammfunktion nachweisen}{\gpunkte[24 23 22]{\textcolor{black!25}{$\circ$}\,\textcolor{black!25}{$\circ$}\,\textcolor{black!80}{$\bullet$}\,\textcolor{black!80}{$\bullet$}\,\textcolor{black!80}{$\bullet$}}}\special{pdf:dest (stufe-stammfunktion-nachweisen) [@thispage /XYZ @xpos @ypos null]}}
% kein Einstieg: keine Schülerfrage in der Gliederung (Schülerfrage, -Fälle, -Form, -Weg, -Skizze; oder Frage) – Einstieg nur von dort, nie aus der Bank (pruefung.md § 5, 12.10.)
% 2024-bebb-gk-B2.1e-m1 (Original 2024-bebb-gk-B2.1e, Schwere I, rang 1; herausgelöste Fassung: keine)
\begin{pfaufg}{GK~’24\\B2.1e}{1.}{}
Gegeben ist die in $\mathbb{R}$ definierte Funktion $f$ mit $f(x) = (x - 3) \cdot e^{-\tfrac{x}{3} + 2}$. Zeige, dass $F$ mit $F(x) = -3x \cdot e^{-\tfrac{x}{3} + 2}$ eine Stammfunktion von $f$ ist.
\karo{2}
\fusshilfe{\mbox{1:~f(x)}}
\end{pfaufg}
% 2023-bebb-gk-B2.1l-m1 (Original 2023-bebb-gk-B2.1l, Schwere II, rang 3; herausgelöste Fassung: keine)
\begin{pfaufg}{GK~’23\\B2.1l}{2.}{}
Gegeben ist die in $\mathbb{R}$ definierte Funktion $f$ mit $f(x) = \tfrac{1}{3}(x^2 - 9) \cdot e^{\tfrac{x}{2}}$; ihr Graph ist $G$. Es gilt $\int \tfrac{1}{3}x^2 \cdot e^{\tfrac{x}{2}}\,dx = \left(\tfrac{2}{3}x^2 - \tfrac{8}{3}x + \tfrac{16}{3}\right) \cdot e^{\tfrac{x}{2}} + C$. Ermittle damit eine Stammfunktion von $f$. Die Abbildung zeigt die Draufsicht einer Tischplatte; ihr Rand ist symmetrisch zu beiden Achsen, im III. Quadranten ist er $G$. 1 LE entspricht 10 cm. Berechne den Flächeninhalt der Tischplatte in Quadratzentimetern.
\gzwei{\begin{abiksys}[xmin=-3.5,xmax=3.5,ymin=-3.5,ymax=3.5,le=0.4]\funktionab[-1.5]{(\x^2-9)*exp(\x/2)/3}{G}{-3}{0}\funktionab{-(\x^2-9)*exp(\x/2)/3}{}{-3}{0}\funktionab{(\x^2-9)*exp(-\x/2)/3}{}{0}{3}\funktionab{-(\x^2-9)*exp(-\x/2)/3}{}{0}{3}\end{abiksys}}{\karo{4}}
\fusshilfe{\mbox{2:~($\tfrac{2}{3}$ x² $-$ $\tfrac{8}{3}$ x $-$ $\tfrac{2}{3}$) · $e^{x/2}$; 1457 cm²}}
\end{pfaufg}
% 2022-bebb-gk-B2.2d-m1 (Original 2022-bebb-gk-B2.2d, Schwere III, rang 2; herausgelöste Fassung: keine)
\begin{pfaufg}{GK~’22\\B2.2d}{3.}{}
Gegeben ist die Funktion $f$ mit $f(x) = -\tfrac{1}{18}x^3 + \tfrac{1}{3}x^2$, $x \in \mathbb{R}$, und die Funktion $F$ mit $F(x) = -\tfrac{1}{72}x^4 + \tfrac{1}{9}x^3$. Weise nach, dass $F$ eine Stammfunktion von $f$ ist. Es gilt: $x = 0$ und $x = 6$ sind die einzigen Nullstellen von $f$, und $F(8) = 0$. Begründe damit: Der Graph von $f$, die $x$-Achse und die Gerade $x = 8$ begrenzen zwei Flächenstücke mit gleichem Inhalt.
\karo{4}
\fusshilfe{\mbox{3:~f(x); gleicher Inhalt}}
\end{pfaufg}

% ===================== Stufe: Integral gesucht =====================
\gkopfvor{\pfstufekopf{Integral gesucht}{\gpunkte[26 25 23]{\textcolor{black!80}{$\bullet$}\,\textcolor{black!80}{$\bullet$}\,\textcolor{black!25}{$\circ$}\,\textcolor{black!80}{$\bullet$}\,\textcolor{black!25}{$\circ$}}}\special{pdf:dest (stufe-integral-gesucht) [@thispage /XYZ @xpos @ypos null]}}
% kein Einstieg: keine Schülerfrage in der Gliederung (Schülerfrage, -Fälle, -Form, -Weg, -Skizze; oder Frage) – Einstieg nur von dort, nie aus der Bank (pruefung.md § 5, 12.10.)
% 2026-bb-gk-A1.1a-m1 (Original 2026-bb-gk-A1.1a, Schwere I, rang 4; herausgelöste Fassung: keine)
\begin{pfaufg}{GK~’26\\A1.1a}{4.}{\pfohnetr}
Gegeben ist die in $\mathbb{R}$ definierte Funktion $f$ mit $f(x) = x^3 + 2x^2 - 3x$; sie hat die Nullstellen $-3$, $0$ und $1$ (Abbildung). Berechne $\int_{-1}^{0} f(x)\,dx$.
\gzwei{\begin{abiksys}[xmin=-3.5,xmax=1.5,ymin=-2,ymax=7,le=0.5]\funktion[1.3]{\x^3+2*\x^2-3*\x}{G_f}\end{abiksys}}{\karo{3}}
\fusshilfe{\mbox{4:~$\tfrac{23}{12}$}}
\end{pfaufg}
% 2025-bebb-gk-B2.2d-m1 (Original 2025-bebb-gk-B2.2d, Schwere II, rang 5; herausgelöste Fassung: keine)
\begin{pfaufg}{GK~’25\\B2.2d}{5.}{}
Gegeben ist die Funktion $f$ mit $f(x) = (3 - x) \cdot e^x$; $F$ mit $F(x) = (4 - x) \cdot e^x$ ist eine Stammfunktion von $f$. Als Näherungswert für $\int_{-2}^{2} f(x)\,dx$ wird $\tfrac{1}{2} \cdot e^2 \cdot 4$ verwendet. Berechne den exakten Wert des Integrals und die prozentuale Abweichung des Näherungswerts davon.
\karo{4}
\fusshilfe{\mbox{5:~13,97; 5,8 \%}}
\end{pfaufg}
% 2023-bebb-gk-B2.2i-m1 (Original 2023-bebb-gk-B2.2i, Schwere II, rang 6; herausgelöste Fassung: keine)
\begin{pfaufg}{GK~’23\\B2.2i}{6.}{}
Die Temperatur in einem Backofen in °C wird in den ersten fünf Minuten nach dem Einschalten beschrieben durch $h(x) = 4x^3 - 45x^2 + 162x + 20$, $0 \le x \le 5$; $x$ ist die Zeit in Minuten. Berechne den Wert des Terms $\tfrac{1}{5} \cdot \int_{0}^{5} h(x)\,dx$ und deute ihn im Sachzusammenhang.
\karo{4}
\fusshilfe{\mbox{6:~175; in den ersten fünf Minuten beträgt die Temperatur im Mittel 175 °C}}
\end{pfaufg}

% ===================== Stufe: Grenze oder Parameter gesucht =====================
\gkopfvor{\pfstufekopf{Grenze oder Parameter gesucht}{\gpunkte[26 24 22]{\textcolor{black!80}{$\bullet$}\,\textcolor{black!25}{$\circ$}\,\textcolor{black!80}{$\bullet$}\,\textcolor{black!25}{$\circ$}\,\textcolor{black!80}{$\bullet$}}}\special{pdf:dest (stufe-grenze-oder-parameter-gesucht) [@thispage /XYZ @xpos @ypos null]}}
% kein Einstieg: keine Schülerfrage in der Gliederung (Schülerfrage, -Fälle, -Form, -Weg, -Skizze; oder Frage) – Einstieg nur von dort, nie aus der Bank (pruefung.md § 5, 12.10.)
% 2024-bebb-gk-B2.2f-m1 (Original 2024-bebb-gk-B2.2f, Schwere II, rang 7; herausgelöste Fassung: keine)
\begin{pfaufg}{GK~’24\\B2.2f}{7.}{}
Gegeben ist die Funktion $f$ mit $f(x) = 0{,}5x^4 - 3x^2 + 2{,}5$, $x \in \mathbb{R}$. Der Graph von $h$ entsteht aus dem Graphen von $f$ durch Stauchung in Richtung der $y$-Achse mit dem Faktor $a$. Der Graph von $h$ und die $x$-Achse begrenzen für $-1 \le x \le 1$ eine Fläche mit dem Inhalt 1. Bestimme $a$.
\karo{4}
\fusshilfe{\mbox{7:~0,31}}
\end{pfaufg}
% 2022-bebb-gk-A1.3a-m1 (Original 2022-bebb-gk-A1.3a, Schwere III, rang 8; herausgelöste Fassung: keine)
\begin{pfaufg}{GK~’22\\A1.3a}{8.}{\pfohnetr}
Gegeben ist die Funktion $f$ mit $f(x) = x^2$. Der Graph von $f$ und die Gerade mit der Gleichung $y = m \cdot x$, $m < 0$, schließen eine Fläche mit dem Inhalt $\tfrac{9}{2}$ ein. Bestimme $m$.
\karo{4}
\fusshilfe{\mbox{8:~Schnittstellen m und 0; $-$3}}
\end{pfaufg}
% 2026-bb-gk-B2.2e-m1 (Original 2026-bb-gk-B2.2e, Schwere III, rang 9; herausgelöste Fassung: keine)
\begin{pfaufg}{GK~’26\\B2.2e}{9.}{}
Gegeben ist die Funktion $f$ mit $f(x) = 2e^x + 3$. Der Graph von $f$, die Gerade $y = 3$, die $y$-Achse und die Gerade $x = u$ mit $u < 0$ begrenzen eine Fläche mit dem Inhalt $2 - 2e^u$. Für $u = -\ln 4$ teilt die Gerade $x = a$ diese Fläche in zwei inhaltsgleiche Teile. Ermittle $a$.
\karo{4}
\fusshilfe{\mbox{9:~$\tfrac{3}{2}$; $-$0,47}}
\end{pfaufg}

\pfblattende{10}{26 & 25 & 24 & 23 & 22 & 21 & 20 & 19 & 18 & 17}{A1.1a B2.2e & B2.2d & B2.1e B2.2f & B2.1l B2.2i & A1.3a B2.2d &  & A1.1a B2.1h & A1.1a B2.2e & B1.2c & B1.2d}%
{\pfausschnitt{}{Integral berechnen}{Flächeninhalt \pkt{} Integral am Graphen}}
\end{document}